\chapter{Discussion}\label{ch:discussion}
The five-date comparison gives mixed results. Some policies reduce estimated carbon on several dates, but none meets every frozen stability condition. Interpreting these changes requires attention to their size, the associated waiting and user burden, and variation between replays. The later V2 tests address a separate question: how candidate selection and feedback from the simulated queue behave in the implementation.

\section{Answers to the research questions}
\subsection{RQ1: Changes in modelled carbon}
The policies did not retain the same ordering across dates. B1 illustrates why a count of positive dates is insufficient: its worst increase outweighed four reductions. B2 also had a negative equal-day mean, while Adaptive Balanced and the day/night hybrid had small positive mean percentage reductions. Table~\ref{tab:final-policy-summary} retains the full numerical comparison.

History-only produced a small positive primary-model mean under several aggregation methods. However, it fell short of the required number of material-positive dates and failed the worst-day condition. The per-job figure divides a shared-queue carbon difference by job count; it cannot be read as an individual's footprint. Its positive mean also remains sensitive to occupancy representation and time anchoring.

For the tested configurations and workloads, RQ1 has a conditional answer: a lower-intensity admission plan can either reduce or increase estimated carbon once Slurm schedules the shared queue. This agrees with the distinction between ideal temporal opportunity and practically constrained shifting discussed by Sukprasert and colleagues~\cite{sukprasert2024limitations}. What matters here is carbon exposure over the resulting execution intervals, not intensity at the selected release timestamp alone.

\subsection{RQ2: Waiting and user burden}
The service question is how much waiting accompanies a carbon change and who bears it. Adaptive's runtime-proportional bound prevented the direct holds beyond configured runtime seen under B1. This protection applied to its chosen holds, not to later queueing. Its mean reported additional total wait was lower than B1's, although the smaller direct-delay budget did not improve every carbon or service outcome.

History-only shows another limitation of the average. On 9 December, a small population mean coexisted with a large share of direct delay concentrated on one user. The cumulative job-hours in Section~\ref{sec:user-burden-results} add the delays across that user's jobs; they do not describe one extremely long hold. Limiting each job's delay and limiting a user's cumulative burden address different service risks.

User burden here means imposed waiting, not a user carbon account. Wassermann and colleagues show that carbon attribution needs explicit choices about measurement and responsibility for idle resources~\cite{wassermann2024usercentric}. The waiting distribution identifies who experiences a service cost, while carbon is assessed for the shared cage scenario.

A carbon reduction need not make service worse for every user. Although direct policy delay is non-negative, total waiting may increase or decrease as jobs compete for resources. Replay variation also contributes to the differences, so each paired change cannot automatically be assigned to deliberate policy action. Positive-only additional waiting reveals later starts that would otherwise be offset by earlier starts elsewhere. Similarly, a favourable population P95 does not mean that every affected job benefited.

\subsection{RQ3: Information and delay safeguards}
Observed runtime helps estimate whether a later interval reduces carbon exposure over a non-preemptive job's complete execution. A completed baseline schedule also reveals future queue and load information. The given-runtime policies and Oracle reference use these advantages for diagnosis. They therefore need to be distinguished from policies that rely on historical predictors.

Earlier-month quantiles remove some of that hindsight from decision runtime and load. Their performance depends on calibration, support and temporal change. History-only delayed 58 evaluation jobs, of which 44 received direct delays longer than their configured execution runtime. A policy can obey a budget based on a prediction while imposing a large delay relative to the eventual job duration. The representative Oracle and History-only rules also have different thresholds and weights, so their comparison does not isolate the effect of prediction error alone.

The Low-impact rule delayed only three evaluation jobs across the five dates, while the Oracle reference made no deliberate delays. This shows when the implemented rules abstained, rather than proving that the workloads had no useful shifting opportunities. The Oracle rule is not a mathematical upper bound. V2 also identified a selection weakness: ranking before checking all constraints can miss a feasible alternative. The regression test demonstrates this code behaviour, but provides no revised results for the frozen policies.

All methods use a historical carbon series without archived issue-time forecasts. The history-only label applies to runtime and load information; it does not describe a complete online forecast evaluation. Temporal-shifting research shows why forecast uncertainty matters, but testing an assumed noise process is different from measuring errors in historical forecasts~\cite{wiesner2021wait}. Runtime, load and carbon information need separate treatment because they come from different sources and have different timing restrictions.

\subsection{RQ4: Consistency across dates and assumptions}
The mean, median, worst date and observed repeat range answer different questions. An equal-day mean weights every selected case equally. The median describes the middle result and is less affected by an extreme date, while the worst date records an adverse outcome that either summary could hide. B1's median reduction was positive at 0.2205\%, even though its mean was negative. Reporting both makes that distribution clear. Standard deviation and interquartile range describe differences across dates; repeat ranges describe variation within a date.

Aggregation changes the question being answered. Adaptive's positive mean daily percentage and negative normalised and pooled savings in Section~\ref{sec:main-results} arise from different weights assigned by daily baseline emissions and job counts. The normalised carbon and waiting means include all evaluation jobs, even those without deliberate delay. Percentiles also require care: averaging daily P95 values does not give a pooled job-level P95, and subtracting two population P95 values does not give the P95 of paired job-level differences.

No policy met the full frozen carbon screen. History-only retained a positive mean direction when any one date was removed, which is useful descriptive evidence. It still failed the material-positive-date and worst-day conditions. The service screen assesses a separate set of costs. Neither screen is a statistical significance test, and every five-date summary retains the negative dates.

Identical-input runs show why repeat variation matters. The 12 no-action cases had no intentional delay to explain their output differences. In particular, Oracle's added waiting cannot be read as a policy-imposed waiting budget. Baseline and Adaptive repeats reveal some variation, but one run per date for each other policy cannot establish its full distribution.

\section{What the carbon model supports}
The framework compares counterfactual schedules under the same power scenario. The trace records resource allocation, but not application activity or the electricity consumed by individual Stanage jobs. Mapping normalised occupancy to the 55 kW load-dependent range is therefore a modelling assumption. The regular-use estimate of 195 kW came without an occupancy measurement. Placing it at full normalised occupancy defines a scenario endpoint; it cannot establish measured full-load power or validate the model's slope.

On the common model-completion horizon, the fixed 140 kW term contributes the same integral to each replay and cancels from the paired absolute difference. It remains in the whole-cage percentage denominator. Extra waiting therefore does not automatically add a different fixed idle-energy charge to the policy. For History-only, the sum of paired savings was approximately 0.3922 kg \co{}, equivalent to a pooled reduction of 0.0813\% in load-dependent carbon but only 0.00335\% in whole-cage carbon. The execution windows overlap, so this sum is not a continuous five-day emissions account or a basis for annualising savings.

This separation follows the distinction between dynamic and infrastructure components in BLOOM~\cite{luccioni2023bloom}, without transferring its component ratios to Stanage.

Disagreement between capacity, node-request and distinct-node views shows sensitivity to the model. History-only's positive primary-model mean was accompanied by negative mean reductions under both alternatives. All views need to remain visible, rather than selecting the favourable one. An aggregate power series could help calibrate the model and identify which occupancy representation best tracks electrical behaviour. For now, the alternatives show how the conclusion changes with the allocation model; they provide no independent electrical measurements.

Trace-based estimation and HPC power modelling likewise separate resource-to-power validity from scheduling correctness~\cite{west2024ichnos,obrien2017survey}. The three occupancy views are retained rather than selecting the most favourable one.

The regional signal provides an accounting proxy for electricity supply, not a measurement of a job's causal effect on generation. The literature separates operational accounting from lifecycle and consequential questions~\cite{lannelongue2021green,gupta2022act,sukprasert2024limitations}. This dissertation reports generation \co{} within the stated cage-power boundary. Legacy \texttt{CO2e} field names do not make it a lifecycle inventory, and the calculation applies no additional PUE multiplier.

In the terminology used by Dodge and colleagues, this is an attributional comparison, while the consequential effect on generation would require marginal-emissions evidence~\cite{dodge2022carbon}. Agreement among occupancy models does not remove that distinction. It checks an allocation assumption inside the experiment, not the causal response of the electricity system.

\section{Replay validity and generalisation}
Matched inputs and common accounting horizons support comparison, but a completed replay can still have timing errors. Section~\ref{sec:timing-audit} reports these errors. The largest late completions occurred after all jobs had started, so they cannot explain earlier waiting; the carbon model also excludes their full overrun. This narrows the impact of that anomaly. Other timing errors may still matter, and the clock-gap trigger remains unresolved.

The final Adaptive repeat also required a retry after 161 jobs timed out in its first attempt. The failed attempt is retained, while the successful attempt supplies the final result. Thus 55 completed replays is the final comparison denominator, not a statement that every attempt succeeded. The older 8 December pairing has a different validity problem: 24 initially running jobs were restored late in the policy run. Its recorded carbon increase and P95 decrease remain historical audit observations, but neither supports a causal explanation about that policy. This invalid pairing is separate from the five-date results.

The later Slurm simulator work documents variability that supports repeating schedules~\cite{simakov2022accurate}. Repetition addresses this source of internal uncertainty, but cannot recover missing production reservations, correct a hardware reconstruction or calibrate a power model. Validation against a physical system requires separate evidence~\cite{jokanovic2018validation}.

Submission-time reconstruction is an additional limitation. Section~\ref{sec:submission-residual-audit} shows that the residual removed by re-anchoring can exceed one half-hour carbon interval, including at the evaluation-job P95 in one baseline. This is not the same error as a late completion event. Section~\ref{sec:negative-wait-audit} also identifies negative waiting markers retained in the original summaries; these are not valid physical waiting intervals. Because jobs can be shifted by different amounts, configured runtime does not by itself preserve logged overlap or establish feasible reconstructed placements. The estimates remain conditional on this accounting convention. The read-only audits do not establish revised savings or rankings, and the maintenance checks do not retroactively validate the original timing.

The five dates were selected for a case study of one cluster; they are not independent random samples of the year. Carry-in work and execution tails can overlap. Each tail drains a finite workload without the full normal arrival stream from later days, which may reduce background competition relative to production. The runs do not establish the direction or size of the resulting policy bias. Assumed flexibility and the fixed basic-priority configuration also limit generalisation. The findings apply to these finite-window cases, without estimating an annual population effect.

The supported completed-work cohort is narrower than the production workload. Failed and cancelled jobs may consume resources without appearing as equivalent replayable jobs here. Runtime caps and carry-in reconstruction also change how source records enter the model. Resource requests and configured runtimes are held fixed between paired policies after this conversion. A complete replay denominator establishes completion of the selected model workload, rather than reproduction of every production job.

\section{What the feedback extension adds}
V2 followed the frozen campaign as a separate engineering extension. It added feasibility checks before candidate ranking and connected a controller to the Slurm simulator event loop. Figure~\ref{fig:v2-feedback} shows the arrived jobs and current queue states available to the controller. It can release held jobs when their expected benefit disappears, while keeping spent waiting charged to the user budget. Unlike a static release plan generated before replay, this controller responds to events during the run. The original seven-policy outcomes remain unchanged.

The three-job functional test shows a limit to cancelling a hold. As Figure~\ref{fig:v2-queue-timeline} shows, the controller cancelled most of the proposed delay, but a larger job had already entered Slurm. The earlier queue position could not be restored. The fixture tests this interaction between admission and queue order; its carbon proxy change does not measure Stanage performance.

The paired 64-job check in Section~\ref{sec:v2-results} used records from 18 November. Both runs started empty, and the controller made no deliberate holds. Their small waiting and proxy differences therefore cannot be attributed to shifting jobs. The single pair cannot separate controller processing from replay variation, and the timing warning remains. This check demonstrates interface operation and rejection behaviour, with no claim about full-day carbon performance.

The 37 saved tests support budget and selection behaviour, not a waiting guarantee from Slurm. V2's carbon proxy also lacks the main study's power scaling and remains separate from its five-date tables.

\section{Practical implications and future work}
The project provides an inspectable comparison of admission rules and their service costs. The results support checking carbon exposure over each proposed execution interval, considering cumulative user burden, and allowing a rule to abstain when it cannot meet its constraints. They also show why direct-delay budgets must be assessed alongside realised queue outcomes. These findings inform the design and evaluation of the tested policies; they do not identify one configuration that is best for every workload.

The frozen comparison and bounded V2 checks are complete. Future work could compare re-anchored and log-aligned accounting on a common horizon, investigate inconsistent event timing, or add metered power, issue-time forecasts, user-approved flexibility and continuous background arrivals. These extensions are outside the completed project. The present findings must be read with the stated limitations, without assuming that those further checks have been performed.
